\section{Introducción}\label{sec:intro}

El Reto AgroCebada FIRA 2026 pide estimar el rendimiento de grano de cebada de temporal, en toneladas por hectárea
(t/ha), de 59 parcelas de Hidalgo, Puebla y Tlaxcala~\cite{bases}. Estas entidades se cuentan entre las principales
productoras de cebada del país~\cite{sader2025}. FIRA busca que la estimación apoye decisiones productivas, de planeación
y eventualmente financieras~\cite{preguntas}.

Contamos con 197 parcelas de nueve municipios, todas del ciclo primavera-verano 2025. En 138 conocemos el rendimiento y
las usamos para entrenar. Las otras 59 son las que se califican, principalmente con la raíz del error cuadrático medio (RMSE), el error absoluto medio y el coeficiente de
determinación, sin ponderar por superficie~\cite{preguntas}. Para todas hay imágenes de los satélites
Sentinel-2~\cite{drusch2012}, Landsat y PlanetScope~\cite{frazier2021}, lluvia~\cite{funk2015},
temperatura~\cite{verdin2020} y topografía.

El rendimiento es el peso en báscula de toda la cosecha, sin ajuste por humedad ni impurezas. Lo reportó sobre todo la
comercializadora que compró el grano y, en menor medida, el productor~\cite{preguntas}. Esto importa porque los
rendimientos reportados suelen traer error~\cite{lobell2020}.

Predecir el rendimiento con satélites es un tema muy estudiado~\cite{vanklompenburg2020,weiss2020}. Aquí el problema
cambia por dos razones. La primera es que 138 parcelas son pocas, y la diferencia entre dos modelos suele ser menor que
el azar de la validación. La segunda es la forma en que se registró el rendimiento.

Al revisar las etiquetas encontramos que cada municipio registra el rendimiento de una sola de dos formas (anexo
técnico, Fig.~\ref{A-fig:regimenes}). En cuatro municipios los valores tienen dos decimales y promedian 4.49\,t/ha. Lo
llamamos régimen A. En los otros cinco están redondeados y promedian 3.05\,t/ha. Lo llamamos régimen B. La diferencia entre ambos, 1.4\,t/ha, es casi tres veces la variación dentro de cada uno,
así que conocer el régimen ya explica buena parte del rendimiento. No sabemos si
el rendimiento se midió distinto o solo se anotó distinto. FIRA pidió tratar todos los registros como una sola variable,
así que usamos el régimen como una característica del municipio y no como dos problemas separados.

Por eso nuestro trabajo no se limita a buscar el algoritmo con menor error. Buscamos responder qué información predice
de verdad el rendimiento, qué tan confiable es el error que estimamos y qué significa ese error para un productor o un
analista.

El resultado es un modelo sencillo. Usa el régimen del municipio y el rendimiento de las parcelas vecinas, y esperamos
que su RMSE en las 59 parcelas sea de 0.46 a 0.49\,t/ha. Lo acompañamos con un
intervalo por parcela, una traducción del error a pesos y una aplicación web.
